\subsection{Amenazas a la validez}

En la discusión se explica cada una de estas amenazas; aquí solo las enumeramos con la forma de evitarlas.

\begin{itemize}
  \item Todo corrió en una sola máquina, en \emph{localhost}, sin red real. Para evitarlo habría que repetir el experimento en un entorno real, con los nodos en máquinas distintas de una red.
  \item 13 de los 4000 envíos de sockets pasaron de 10~ms (el máximo fue 27,6~ms) y no identificamos la causa. Ensanchan el intervalo de confianza de los mensajes por segundo. Lo evitaría cerrando procesos pesados en segundo plano durante la prueba, fijando la afinidad y prioridad de CPU de los nodos en el sistema operativo para evitar cambios de contexto bruscos, y haciendo más repeticiones para amortiguar picos aislados del planificador.
  \item 8 de las 30 corridas de Lamport tardaron cerca de 0,5~s en arrancar, porque un nodo intentó conectarse antes de que el otro abriera su puerto (sección~\ref{sec:discusion-lamport}). Para evitarlo, como se propone en la discusión, cada nodo debería esperar a que todos los servidores estén escuchando antes del primer envío. Ese cambio iría en \texttt{src/comun/nodo.py}.
  \item La ronda de dos fichas da 2, 6 y 8 intercambios, y la hipótesis esperaba 2, 4 y 6. Lo evitaría fijando primero el diseño exacto de la topología (si la ronda corre con una o con dos fichas simultáneas) antes de redactar la hipótesis, o calculando la predicción teórica de H4 directamente sobre el flujo real de mensajes de las dos fichas antes de ejecutar las corridas.
\end{itemize}
